\problemname{BrainFsckVM}
Given a brainfsck program, determine whether it terminates or enters
an endless loop.\\

A brainfsck interpreter has a data array (consisting of unsigned 8-bit
integers) with an index, the so called ``\emph{data index}''; the
entry of the array pointed to by the data index is the so called
``\emph{current entry}''. A brainfsck program consists of a sequence
of the following eight instructions:

\begin{center}
  \begin{tabular}{|l|l|}
    \hline
    \verb?-? & decrease \emph{current entry} by $1$ (modulo $2^8$)\\
    \hline
    \verb?+? & increase \emph{current entry} by $1$ (modulo $2^8$)\\
    \hline
    \verb+<+ & decrease \emph{data index} by $1$ \\
    \hline
    \verb+>+ & increase \emph{data index} by $1$ \\
    \hline
    \verb+[+ & jump behind the matching \verb+]+, 
    if the \emph{current entry} is equal to 0 \\
    \hline
    \verb+]+ & jump behind the matching \verb+[+ if 
    the \emph{current entry} is \emph{not} equal to $0$\\
    \hline
    \verb+.+ & print the \emph{current entry} as character\\
    \hline
    \verb+,+ & read a character and save it into the \emph{current
    entry}. On end of input save $255$.\\
    \hline
  \end{tabular}
\end{center}

Interpretation of a brainfsck program starts at the first instruction,
and terminates if the instruction pointer leaves the end of the
program. After an instruction is executed, the instruction pointer
advances to the instruction to the right (except, of course, if the
loop instructions \verb+[+ or \verb+]+ cause a jump).\\

In addition to the program, you will be given the size of the data
array. The entries of the data array shall be unsigned 8-bit integers,
with usual over- or underflow behaviour. 
At the start of the program, the data array entries and the data 
index shall be initialized to zero.\\

Incrementing or decrementing the data index beyond the
boundaries of the data array shall make it re-enter the data array at
the other boundary; e.g.\ decrementing the data index when it is
zero shall set it to the size of the data array minus one.\\

The \verb+[+ and \verb+]+ instructions define loops and are allowed to
nest. Every given program will be well-formed, i.e. when traversing
the program from left to right, the number of \verb+[+ instructions
minus the number of \verb+]+ instructions will always be greater or
equal zero, and at the end of the program it will be equal to zero.\\

For the purposes of the problem, discard the output of the interpreted program.
\footnote{For purposes of debugging, you may examine it. The third program
	from the sample input would print ``ICPC''; the fourth program,
	when given an arbitrary string, will print a brainfsck program, that in
	turn will print said arbitrary string when executed.}



\section*{Input}
The input starts with a line containing the number of test cases $t$ ($0 < t \le 20$). Each
test case consists of three lines. The first line contains three integers $s_m,
s_c, s_i$, describing the size of the memory, the size of the program code and
the size of the input ($0 < s_m \le 100\,000; 0 < s_c, s_i < 4\,096$).

The second line contains the brainfsck program code to be executed; it
consists of $s_c$ characters.

The third line contains the input of the program, as text (only printable,
non-whitespace characters). Once the program has consumed all available input,
the input instruction shall set the current cell to $255$.

\section*{Output}
For each test case, print one line, containing either ``{\tt Terminates}'' or
``{\tt Loops}'', depending on whether the program either terminates after a
finite number of steps, or enters an endless loop. If the program loops, also
print the indices ($0$-based) of the two \verb+[+ and the \verb+]+ symbols in
the code array that correspond to the endless loop. You may safely assume that
after $50\,000\,000$ instructions, a program either terminated or hangs in an
endless loop, which then was executed at least once. During each iteration of the endless loop
 at most $50\,000\,000$ instructions are executed.
